% Part: normal-modal-logic
% Chapter: filtrations
% Section: euclidean-filtrations

\documentclass[../../../include/open-logic-section]{subfiles}

\begin{document}

\olfileid{nml}{fil}{euc}

\olsection{ইউক্লিডীয় মডেলের ফিল্ট্রেশন}

% BN-SRC-363: ভাগ-মডেলের প্রয়োজনীয় তিরগুলি সমতুল্যতা-শ্রেণির মধ্যে।
\olref[acc]{sec}-এর পদ্ধতি কেবল ইউক্লিডীয়, অথবা সিরিয়াল ও ইউক্লিডীয় মডেলের ক্ষেত্রে কাজ করে না। \olref{fig:ser-eucl}-এর ওপরের মডেলটি দেখো; এটি ইউক্লিডীয় এবং সিরিয়াল উভয়ই। ধরা যাক $\Gamma = \{p, \Box p \}$। $\Gamma$-র মধ্য দিয়ে ফিল্ট্রেশন নিলে $[w_1] = [w_3]$, কারণ $w_1$ ও $w_3$-ই একমাত্র জগৎ যারা~$\Gamma$-তে একমত। যেকোনো ফিল্ট্রেশনেও $\mModel{M}$ থেকে উত্তরাধিকারসূত্রে পাওয়া তিরগুলি থাকবে, যেমন \olref{fig:ser-eucl2}-এ। ওই ভাগ-মডেল ইউক্লিডীয় নয়।
% BN-SRC-724 TARGET-CORRECTION-BEGIN
\emph{উৎস-সংশোধন:} মূল ছবিতে দ্বিতীয় জগতের আত্মতির বাদ পড়েছে। প্রথম জগত থেকে সেখানে তির থাকায় ইউক্লিডীয়তা ওই আত্মতির বাধ্য করে; সিরিয়ালিটির জন্যও সেটি দরকার। আদি ও ভাগ-মডেল উভয় ছবিতে দ্বিতীয় জগতের আত্মতির যোগ হয়েছে। ওই জগতে চলটি সত্য বলে আবশ্যিক সূত্রের প্রদর্শিত সত্যমান বজায় থাকে।
% BN-SRC-724 TARGET-CORRECTION-END
 তাছাড়া ওই মডেলে তির যোগ করেও তাকে ইউক্লিডীয় করা যায় না। $[w_2]$ ও $[w_4]$-এর মধ্যে উভমুখী তির যোগ করতে হতো, তারপর $[w_2]$ ও~$[w_5]$-এর মধ্যেও। কিন্তু~$w_2$-তে $\Box
p$ সত্য থাকার কথা, অথচ~$w_5$-তে $p$ মিথ্যা।

\begin{figure}[htpb]
  \centering
  \begin{tikzpicture}[modal]
    \node[world] (w1)
          [label={left: \mFalse{p}},
            label={below: $\mSat{{}}{\Box p}$}] {$w_1$};
    \node[world] (w2)
          [label={right: \mTrue{p}},
            label = {below: $\mSat{{}}{\Box p}$},
            right=of w1] {$w_2$};
    \draw[->] (w1) to (w2);
    \draw[reflexive above] (w2) to (w2);
    \node[world] (w3)
          [label={left: \mFalse{p}},
            label={below: $\mSat{{}}{\Box p}$},
            below=of w1] {$w_3$};
    \node[world] (w4)
          [label={right:\mTrue{p}},
            label={below: $\mSat/{{}}{\Box p}$},
            right=of w3] {$w_4$};
    \draw[->] (w3) to (w4);
    \draw[reflexive above] (w4) to (w4);
    \node[world] (w5)
          [label={right: \mFalse{p}},
            label=below:$\mSat/{{}}{\Box p}$,
            right=of w4] {$w_5$};
    \draw[->, bend left] (w4) to (w5);
    \draw[->, bend left] (w5) to (w4);
    \draw[reflexive above] (w5) to (w5);
  \end{tikzpicture}
  \caption{একটি সিরিয়াল ও ইউক্লিডীয় মডেল।}
  \ollabel{fig:ser-eucl}
\end{figure}

\begin{figure}[ht]
  \centering
  \begin{tikzpicture}[modal,every text node part/.style={align=center}]
    \node[world] (w1)
          [label={left: \mFalse{p}},
            label={right:$[w_1]=[w_3]$},
            label={below: $\mSat{{}}{\Box p}$}] {$[w_1]$};
    \node[world] (w2)
         [label={right: \mTrue{p}},
           label = {below: $\mSat{{}}{\Box p}$},
           right=of w1, yshift=1.5cm] {$[w_2]$};
    \draw[->] (w1) to (w2);
    \draw[reflexive above] (w2) to (w2);
    \node[world] (w4)
         [label={right:\mTrue{p}},
           label={below: $\mSat/{{}}{\Box p}$},
           right=of w1,yshift=-1.5cm] {$[w_4]$};
    \draw[->] (w1) to (w4);
    \draw[reflexive above] (w4) to (w4);
    \node[world] (w5)
         [label={right: \mFalse{p}},
             label=below:$\mSat/{{}}{\Box p}$,
             right=of w4] {$[w_5]$};
    \draw[->, bend left] (w4) to (w5);
    \draw[->, bend left] (w5) to (w4);
    \draw[reflexive above] (w5) to (w5);
  \end{tikzpicture}
  \caption{\olref{fig:ser-eucl}-এর মডেলের ফিল্ট্রেশন।}
  \ollabel{fig:ser-eucl2}
\end{figure}

বিশেষত, ইউক্লিডীয় ফিল্ট্রেশন পেতে উপ-!!{formula}s-এর অধীনে বদ্ধ নির্বিচার $\Gamma$-র মধ্য দিয়ে ফিল্ট্রেশন নিলেই চলে না। বরং \emph{মোডালভাবে বদ্ধ} সমষ্টি $\Gamma$ নিতে হয় (\olref[pre]{defn:modallyclosed} দেখুন)। এমন অশূন্য বাক্যসমষ্টি অসীম; তাই এগুলি থেকে সরাসরি সসীম মডেল ধর্ম বা সংশ্লিষ্ট পদ্ধতির সিদ্ধান্তযোগ্যতা পাওয়া যায় না।
% BN-SRC-725 TARGET-CORRECTION-BEGIN
\emph{উৎস-সংশোধন:} খালি সমষ্টিও মোডালভাবে বদ্ধ। অসীমতার দাবিটি অশূন্য সমষ্টির জন্য: একটি সূত্র থাকলে তার সামনে বারবার আবশ্যিকতার চিহ্ন যোগ করে অনন্ত ভিন্ন সূত্র পাওয়া যায়। অসীম সূত্রসমষ্টি থেকে ভাগ-মডেল সব ক্ষেত্রেই অসীম, এমন দাবি করা হচ্ছে না।
% BN-SRC-725 TARGET-CORRECTION-END


\begin{thm}\ollabel{thm:modal-closed-filt}
  ধরা যাক $\Gamma$ মোডালভাবে বদ্ধ, $\mModel{M}=\tuple{W,R,V}$ এবং $\mModel{M^*} = \tuple{W^*,R^*,V^*}$ হলো এই মোডালভাবে বদ্ধ সমষ্টির মধ্য দিয়ে $\mModel{M}$-এর সবচেয়ে স্থূল ফিল্ট্রেশন।
  \begin{enumerate}
  \item $\mModel{M}$ প্রতিসম হলে $\mModel{M^*}$-ও প্রতিসম।
  \item $\mModel{M}$ পরিযায়ী হলে $\mModel{M^*}$-ও পরিযায়ী।
  \item $\mModel{M}$ ইউক্লিডীয় হলে $\mModel{M^*}$-ও ইউক্লিডীয়।
  \end{enumerate}
\end{thm}

\begin{proof}
% BN-SRC-364: প্রমাণের ক্ষেত্রগুলি উপপাদ্যের প্রতিসাম্য, পরিযায়িতা ও ইউক্লিডীয়তার ক্রমে সাজানো।
\begin{enumerate}
  \item অনুশীলনী। সব প্রতিসম মডেলে \Ax{B} ও $\Ax{B_\Diamond}$ বৈধ, এই তথ্য ব্যবহার করুন।
  \item $\mModel{M^*}$ সবচেয়ে স্থূল ফিল্ট্রেশন হলে সংজ্ঞা অনুসারে $R^*[u][v]$ তখনই এবং কেবল তখনই এবং কেবল তখনই, যখন $C_1(u,v)$। পরিযায়িতা প্রমাণে ধরি $C_1(u,v)$ ও $C_1(v,w)$; দেখাতে হবে $C_1(u,w)$। \iftag{prvBox}{ধরি $\mSat{M}{\Box !A}[u]$। সব পরিযায়ী মডেলে \Ax{4} বৈধ বলে $\mSat{M}{\Box\Box!A}[u]$; বদ্ধতার কারণে $\Box\Box!A \in \Gamma$, তাই $C_1(u,v)$ থেকে $\mSat{M}{\Box!A}[v]$ এবং $C_1(v,w)$ থেকে $\mSat{M}{!A}[w]$। }{}\iftag{prvDiamond}{ধরি $\mSat{M}{!A}[w]$। মোডাল বদ্ধতার কারণে $\Diamond !A \in \Gamma$, তাই $C_1(v,w)$ থেকে $\mSat{M}{\Diamond !A}[v]$। আবার মোডাল বদ্ধতায় $\Diamond\Diamond!A \in \Gamma$; $C_1(u,v)$ থেকে $\mSat{M}{\Diamond\Diamond !A}[u]$। সব পরিযায়ী মডেলে $\Ax{4}_\Diamond$ বৈধ বলে $\mSat{M}{\Diamond!A}[u]$।}{}
  \item অনুশীলনী। সব ইউক্লিডীয় মডেলে \Ax{5} ও $\Ax{5_\Diamond}$ বৈধ, এই তথ্য ব্যবহার করুন।
\end{enumerate}
\end{proof}

\begin{prob}
  \olref[nml][fil][euc]{thm:modal-closed-filt}-এর প্রমাণ সম্পূর্ণ করুন।
\end{prob}

\end{document}
